%%%%%%%%%%%%%%%%%%%%%%%%%%%%%%%%%%%%%%%%%%%%%%%%%%%%%%%%%%%%
% 6.12 MATROID THEORY
% The Composition of Advanced Mathematics
%%%%%%%%%%%%%%%%%%%%%%%%%%%%%%%%%%%%%%%%%%%%%%%%%%%%%%%%%%%%

\section{Matroid Theory}
\label{sec:matroid-theory}

\prerequisite{
Sets, functions, graph theory, trees, spanning trees, linear
independence, vector spaces, finite fields, basic combinatorics,
optimization, and familiarity with rank and dimension.
}

\learningobjectives{
By the end of this section, the reader should be able to:
\begin{itemize}
    \item explain the purpose of matroid theory;
    \item define a matroid using independent sets;
    \item verify the hereditary and exchange axioms;
    \item identify dependent sets, bases, circuits, and rank;
    \item understand graphic and vector matroids;
    \item work with uniform and partition matroids;
    \item recognize spanning sets, closure, and flats;
    \item define the rank function of a matroid;
    \item use rank axioms and submodularity;
    \item understand basis exchange;
    \item define matroid duality;
    \item identify cocircuits and cobases;
    \item perform deletion and contraction;
    \item understand matroid minors;
    \item recognize representable matroids;
    \item understand binary matroids;
    \item connect matroids with projective geometry;
    \item understand the greedy algorithm characterization of matroids;
    \item formulate maximum-weight basis problems;
    \item recognize the Tutte polynomial of a matroid;
    \item understand direct sums and connectedness;
    \item recognize matroid intersection and union problems;
    \item connect matroid theory with graphs, linear algebra, geometry,
          optimization, coding theory, and design theory.
\end{itemize}
}

%%%%%%%%%%%%%%%%%%%%%%%%%%%%%%%%%%%%%%%%%%%%%%%%%%%%%%%%%%%%
\subsection{Why Matroids?}
\label{subsec:why-matroids}
%%%%%%%%%%%%%%%%%%%%%%%%%%%%%%%%%%%%%%%%%%%%%%%%%%%%%%%%%%%%

Many different areas of mathematics contain a notion of
\emph{independence}.

In linear algebra, a set of vectors may be linearly independent.

In graph theory, a set of edges may contain no cycle.

In combinatorics, one may choose objects subject to compatibility
constraints.

These ideas look different, but they satisfy strikingly similar
structural rules.

Matroid theory abstracts these common properties.

The central idea is

\[
\boxed{
\text{matroid}
=
\text{abstract theory of independence}.
}
\]

%%%%%%%%%%%%%%%%%%%%%%%%%%%%%%%%%%%%%%%%%%%%%%%%%%%%%%%%%%%%
\subsection{Independent-Set Definition}
\label{subsec:matroid-independent-definition}
%%%%%%%%%%%%%%%%%%%%%%%%%%%%%%%%%%%%%%%%%%%%%%%%%%%%%%%%%%%%

\begin{definitionbox}{Matroid}
A matroid is a pair

\[
\boxed{
M=(E,\mathcal{I}),
}
\]

where \(E\) is a finite ground set and

\[
\mathcal{I}\subseteq\mathcal{P}(E)
\]

is a family of subsets called independent sets satisfying:

\begin{enumerate}
    \item
    \[
    \varnothing\in\mathcal{I};
    \]

    \item if
    \[
    I\in\mathcal{I}
    \]
    and
    \[
    J\subseteq I,
    \]
    then
    \[
    J\in\mathcal{I};
    \]

    \item if
    \[
    I,J\in\mathcal{I}
    \]
    and
    \[
    |I|<|J|,
    \]
    then there exists
    \[
    e\in J\setminus I
    \]
    such that
    \[
    I\cup\{e\}\in\mathcal{I}.
    \]
\end{enumerate}
\end{definitionbox}

The third condition is called the exchange axiom.

%%%%%%%%%%%%%%%%%%%%%%%%%%%%%%%%%%%%%%%%%%%%%%%%%%%%%%%%%%%%
\subsection{The Ground Set}
\label{subsec:matroid-ground-set}
%%%%%%%%%%%%%%%%%%%%%%%%%%%%%%%%%%%%%%%%%%%%%%%%%%%%%%%%%%%%

The set

\[
E
\]

is called the ground set.

Its elements depend on the application.

For example:

\begin{itemize}
    \item in a graph matroid, \(E\) is a set of graph edges;
    \item in a vector matroid, \(E\) indexes vectors;
    \item in a partition matroid, \(E\) is divided into categories.
\end{itemize}

The matroid structure determines which subsets of \(E\) are considered
independent.

%%%%%%%%%%%%%%%%%%%%%%%%%%%%%%%%%%%%%%%%%%%%%%%%%%%%%%%%%%%%
\subsection{The Hereditary Axiom}
\label{subsec:matroid-hereditary}
%%%%%%%%%%%%%%%%%%%%%%%%%%%%%%%%%%%%%%%%%%%%%%%%%%%%%%%%%%%%

The hereditary axiom states:

\[
\boxed{
I\in\mathcal{I}
\text{ and }
J\subseteq I
\Longrightarrow
J\in\mathcal{I}.
}
\]

Thus removing elements from an independent set cannot create
dependence.

This matches both:

\[
\text{linear independence}
\]

and

\[
\text{acyclic edge sets}.
\]

%%%%%%%%%%%%%%%%%%%%%%%%%%%%%%%%%%%%%%%%%%%%%%%%%%%%%%%%%%%%
\subsection{The Exchange Axiom}
\label{subsec:matroid-exchange}
%%%%%%%%%%%%%%%%%%%%%%%%%%%%%%%%%%%%%%%%%%%%%%%%%%%%%%%%%%%%

Suppose

\[
I,J\in\mathcal{I}
\]

with

\[
|I|<|J|.
\]

The exchange axiom guarantees that some element of \(J\) can be added
to \(I\) while preserving independence.

Thus a smaller independent set can always be enlarged toward a larger
one.

This axiom produces the remarkable uniformity of maximal independent
sets.

%%%%%%%%%%%%%%%%%%%%%%%%%%%%%%%%%%%%%%%%%%%%%%%%%%%%%%%%%%%%
\subsection{Dependent Sets}
\label{subsec:matroid-dependent-sets}
%%%%%%%%%%%%%%%%%%%%%%%%%%%%%%%%%%%%%%%%%%%%%%%%%%%%%%%%%%%%

A subset

\[
D\subseteq E
\]

is dependent if

\[
D\notin\mathcal{I}.
\]

By the hereditary axiom, every superset of a dependent set is also
dependent.

Thus dependence propagates upward under inclusion.

%%%%%%%%%%%%%%%%%%%%%%%%%%%%%%%%%%%%%%%%%%%%%%%%%%%%%%%%%%%%
\subsection{Bases}
\label{subsec:matroid-bases}
%%%%%%%%%%%%%%%%%%%%%%%%%%%%%%%%%%%%%%%%%%%%%%%%%%%%%%%%%%%%

\begin{definitionbox}{Basis of a Matroid}
A basis is a maximal independent subset of \(E\).
\end{definitionbox}

Thus \(B\subseteq E\) is a basis when:

\[
B\in\mathcal{I}
\]

but

\[
B\cup\{e\}\notin\mathcal{I}
\]

for every

\[
e\in E\setminus B.
\]

%%%%%%%%%%%%%%%%%%%%%%%%%%%%%%%%%%%%%%%%%%%%%%%%%%%%%%%%%%%%
\subsection{All Bases Have the Same Size}
\label{subsec:all-bases-same-size}
%%%%%%%%%%%%%%%%%%%%%%%%%%%%%%%%%%%%%%%%%%%%%%%%%%%%%%%%%%%%

\begin{theorem}
Every basis of a finite matroid has the same cardinality.
\end{theorem}

Suppose \(B_1\) and \(B_2\) are bases.

If

\[
|B_1|<|B_2|,
\]

the exchange axiom would allow some

\[
e\in B_2\setminus B_1
\]

to be added to \(B_1\) while preserving independence.

This contradicts maximality of \(B_1\).

Therefore,

\[
\boxed{
|B_1|=|B_2|.
}
\]

%%%%%%%%%%%%%%%%%%%%%%%%%%%%%%%%%%%%%%%%%%%%%%%%%%%%%%%%%%%%
\subsection{Rank of a Matroid}
\label{subsec:matroid-rank}
%%%%%%%%%%%%%%%%%%%%%%%%%%%%%%%%%%%%%%%%%%%%%%%%%%%%%%%%%%%%

\begin{definitionbox}{Rank}
The rank of a matroid \(M\), written

\[
\boxed{
r(M),
}
\]

is the size of any basis of \(M\).
\end{definitionbox}

Since all bases have equal size, this is well-defined.

%%%%%%%%%%%%%%%%%%%%%%%%%%%%%%%%%%%%%%%%%%%%%%%%%%%%%%%%%%%%
\subsection{Rank of a Subset}
\label{subsec:rank-of-subset}
%%%%%%%%%%%%%%%%%%%%%%%%%%%%%%%%%%%%%%%%%%%%%%%%%%%%%%%%%%%%

For

\[
A\subseteq E,
\]

define

\[
\boxed{
r(A)
=
\max
\{
|I|:
I\subseteq A,\;
I\in\mathcal{I}
\}.
}
\]

Thus \(r(A)\) is the size of the largest independent subset contained
in \(A\).

In particular,

\[
\boxed{
r(E)=r(M).
}
\]

%%%%%%%%%%%%%%%%%%%%%%%%%%%%%%%%%%%%%%%%%%%%%%%%%%%%%%%%%%%%
\subsection{Rank Properties}
\label{subsec:rank-properties-matroid}
%%%%%%%%%%%%%%%%%%%%%%%%%%%%%%%%%%%%%%%%%%%%%%%%%%%%%%%%%%%%

The rank function satisfies:

\[
\boxed{
0\leq r(A)\leq|A|;
}
\]

if

\[
A\subseteq B,
\]

then

\[
\boxed{
r(A)\leq r(B);
}
\]

and

\[
\boxed{
r(A\cup B)+r(A\cap B)
\leq
r(A)+r(B).
}
\]

The final inequality is called submodularity.

%%%%%%%%%%%%%%%%%%%%%%%%%%%%%%%%%%%%%%%%%%%%%%%%%%%%%%%%%%%%
\subsection{Submodularity}
\label{subsec:matroid-submodularity}
%%%%%%%%%%%%%%%%%%%%%%%%%%%%%%%%%%%%%%%%%%%%%%%%%%%%%%%%%%%%

A set function \(f\) is submodular if

\[
\boxed{
f(A)+f(B)
\geq
f(A\cup B)+f(A\cap B).
}
\]

Submodularity may be interpreted as a discrete form of diminishing
returns.

Matroid rank is one of the most important examples of a submodular
function.

%%%%%%%%%%%%%%%%%%%%%%%%%%%%%%%%%%%%%%%%%%%%%%%%%%%%%%%%%%%%
\subsection{Circuits}
\label{subsec:matroid-circuits}
%%%%%%%%%%%%%%%%%%%%%%%%%%%%%%%%%%%%%%%%%%%%%%%%%%%%%%%%%%%%

\begin{definitionbox}{Circuit}
A circuit is a minimal dependent set.
\end{definitionbox}

Thus \(C\subseteq E\) is a circuit if:

\[
C
\]

is dependent, but every proper subset of \(C\) is independent.

Circuits represent the smallest possible causes of dependence.

%%%%%%%%%%%%%%%%%%%%%%%%%%%%%%%%%%%%%%%%%%%%%%%%%%%%%%%%%%%%
\subsection{Circuit Characterization}
\label{subsec:circuit-characterization}
%%%%%%%%%%%%%%%%%%%%%%%%%%%%%%%%%%%%%%%%%%%%%%%%%%%%%%%%%%%%

A set \(I\subseteq E\) is independent if and only if it contains no
circuit.

Therefore,

\[
\boxed{
I\in\mathcal{I}
\iff
\text{\(I\) contains no circuit}.
}
\]

This is analogous to graph theory, where an edge set is acyclic exactly
when it contains no cycle.

%%%%%%%%%%%%%%%%%%%%%%%%%%%%%%%%%%%%%%%%%%%%%%%%%%%%%%%%%%%%
\subsection{Circuit Elimination}
\label{subsec:circuit-elimination}
%%%%%%%%%%%%%%%%%%%%%%%%%%%%%%%%%%%%%%%%%%%%%%%%%%%%%%%%%%%%

Circuits satisfy an important elimination property.

If \(C_1\) and \(C_2\) are distinct circuits and

\[
e\in C_1\cap C_2,
\]

then there exists a circuit \(C_3\) satisfying

\[
\boxed{
C_3
\subseteq
(C_1\cup C_2)\setminus\{e\}.
}
\]

This is called the circuit elimination axiom.

It provides an alternative axiomatization of matroids.

%%%%%%%%%%%%%%%%%%%%%%%%%%%%%%%%%%%%%%%%%%%%%%%%%%%%%%%%%%%%
\subsection{Graphic Matroids}
\label{subsec:graphic-matroids}
%%%%%%%%%%%%%%%%%%%%%%%%%%%%%%%%%%%%%%%%%%%%%%%%%%%%%%%%%%%%

One of the most important examples comes from graphs.

Let

\[
G=(V,E)
\]

be a finite graph.

Define a subset of edges to be independent when it contains no cycle.

Then

\[
\boxed{
M(G)
=
(E,\mathcal{I})
}
\]

is a matroid called the graphic matroid or cycle matroid of \(G\).

%%%%%%%%%%%%%%%%%%%%%%%%%%%%%%%%%%%%%%%%%%%%%%%%%%%%%%%%%%%%
\subsection{Independent Sets in Graphic Matroids}
\label{subsec:graphic-independent-sets}
%%%%%%%%%%%%%%%%%%%%%%%%%%%%%%%%%%%%%%%%%%%%%%%%%%%%%%%%%%%%

For a graphic matroid:

\[
\boxed{
I\subseteq E
\text{ is independent}
\iff
I\text{ is a forest}.
}
\]

Thus forests are exactly the independent sets.

If \(G\) is connected, the bases are spanning trees.

%%%%%%%%%%%%%%%%%%%%%%%%%%%%%%%%%%%%%%%%%%%%%%%%%%%%%%%%%%%%
\subsection{Bases of a Graphic Matroid}
\label{subsec:graphic-matroid-bases}
%%%%%%%%%%%%%%%%%%%%%%%%%%%%%%%%%%%%%%%%%%%%%%%%%%%%%%%%%%%%

Suppose \(G\) is connected.

A maximal acyclic edge set is a spanning tree.

Therefore,

\[
\boxed{
\text{bases of }M(G)
=
\text{spanning trees of }G.
}
\]

Since every spanning tree has

\[
|V|-1
\]

edges,

\[
\boxed{
r(M(G))
=
|V|-1.
}
\]

%%%%%%%%%%%%%%%%%%%%%%%%%%%%%%%%%%%%%%%%%%%%%%%%%%%%%%%%%%%%
\subsection{Graphic Matroid Rank}
\label{subsec:graphic-matroid-rank}
%%%%%%%%%%%%%%%%%%%%%%%%%%%%%%%%%%%%%%%%%%%%%%%%%%%%%%%%%%%%

For any edge set

\[
A\subseteq E,
\]

let

\[
c(A)
\]

be the number of connected components of the spanning subgraph

\[
(V,A).
\]

Then

\[
\boxed{
r(A)
=
|V|-c(A).
}
\]

For the entire graph,

\[
\boxed{
r(E)
=
|V|-c(G).
}
\]

%%%%%%%%%%%%%%%%%%%%%%%%%%%%%%%%%%%%%%%%%%%%%%%%%%%%%%%%%%%%
\subsection{Circuits of a Graphic Matroid}
\label{subsec:graphic-matroid-circuits}
%%%%%%%%%%%%%%%%%%%%%%%%%%%%%%%%%%%%%%%%%%%%%%%%%%%%%%%%%%%%

The minimal dependent edge sets in a graph are cycles.

Therefore,

\[
\boxed{
\text{circuits of }M(G)
=
\text{cycles of }G.
}
\]

Thus the abstract word ``circuit'' was chosen to generalize graph
cycles.

%%%%%%%%%%%%%%%%%%%%%%%%%%%%%%%%%%%%%%%%%%%%%%%%%%%%%%%%%%%%
\subsection{Worked Example: Graphic Matroid}
\label{subsec:worked-graphic-matroid}
%%%%%%%%%%%%%%%%%%%%%%%%%%%%%%%%%%%%%%%%%%%%%%%%%%%%%%%%%%%%

\begin{workedexamplebox}{A Cycle Graph}

Consider

\[
G=C_4.
\]

Its edge set contains four edges.

Any set of at most three edges is acyclic.

The entire four-edge set forms a cycle.

Therefore the unique circuit is

\[
E(C_4).
\]

Any three-edge subset is a spanning tree and hence a basis.

\begin{answerbox}
\[
\boxed{
r(M(C_4))=3.
}
\]
\end{answerbox}

\end{workedexamplebox}

%%%%%%%%%%%%%%%%%%%%%%%%%%%%%%%%%%%%%%%%%%%%%%%%%%%%%%%%%%%%
\subsection{Vector Matroids}
\label{subsec:vector-matroids}
%%%%%%%%%%%%%%%%%%%%%%%%%%%%%%%%%%%%%%%%%%%%%%%%%%%%%%%%%%%%

Let

\[
v_1,\ldots,v_n
\]

be vectors over a field \(\mathbb{F}\).

Take

\[
E=\{1,\ldots,n\}.
\]

Declare

\[
I\subseteq E
\]

independent when

\[
\{v_i:i\in I\}
\]

is linearly independent.

This defines a matroid.

%%%%%%%%%%%%%%%%%%%%%%%%%%%%%%%%%%%%%%%%%%%%%%%%%%%%%%%%%%%%
\subsection{Linear Matroid}
\label{subsec:linear-matroid}
%%%%%%%%%%%%%%%%%%%%%%%%%%%%%%%%%%%%%%%%%%%%%%%%%%%%%%%%%%%%

\begin{definitionbox}{Vector Matroid}
The matroid defined by linear independence among a finite collection of
vectors is called a vector matroid or linear matroid.
\end{definitionbox}

The exchange axiom follows from the familiar basis-exchange property of
vector spaces.

%%%%%%%%%%%%%%%%%%%%%%%%%%%%%%%%%%%%%%%%%%%%%%%%%%%%%%%%%%%%
\subsection{Matrix Representation}
\label{subsec:matroid-matrix-representation}
%%%%%%%%%%%%%%%%%%%%%%%%%%%%%%%%%%%%%%%%%%%%%%%%%%%%%%%%%%%%

Let

\[
A
\]

be a matrix over a field \(\mathbb{F}\).

Label its columns by the elements of \(E\).

A set of labels is independent when the corresponding columns are
linearly independent.

The resulting matroid is denoted informally by

\[
\boxed{
M[A].
}
\]

Different matrices can represent the same matroid.

%%%%%%%%%%%%%%%%%%%%%%%%%%%%%%%%%%%%%%%%%%%%%%%%%%%%%%%%%%%%
\subsection{Worked Example: Vector Matroid}
\label{subsec:worked-vector-matroid}
%%%%%%%%%%%%%%%%%%%%%%%%%%%%%%%%%%%%%%%%%%%%%%%%%%%%%%%%%%%%

\begin{workedexamplebox}{Three Vectors in \(\mathbb{R}^2\)}

Consider

\[
v_1=
\begin{pmatrix}
1\\
0
\end{pmatrix},
\qquad
v_2=
\begin{pmatrix}
0\\
1
\end{pmatrix},
\qquad
v_3=
\begin{pmatrix}
1\\
1
\end{pmatrix}.
\]

Every pair is linearly independent.

However,

\[
v_3=v_1+v_2.
\]

Therefore,

\[
\{1,2,3\}
\]

is dependent.

It is minimal dependent, so it is a circuit.

\begin{answerbox}
\[
\boxed{
r(M)=2,
\qquad
\{1,2,3\}\text{ is a circuit}.
}
\]
\end{answerbox}

\end{workedexamplebox}

%%%%%%%%%%%%%%%%%%%%%%%%%%%%%%%%%%%%%%%%%%%%%%%%%%%%%%%%%%%%
\subsection{Representable Matroids}
\label{subsec:representable-matroids}
%%%%%%%%%%%%%%%%%%%%%%%%%%%%%%%%%%%%%%%%%%%%%%%%%%%%%%%%%%%%

\begin{definitionbox}{Representable Matroid}
A matroid is representable over a field \(\mathbb{F}\) if it is
isomorphic to a vector matroid represented by a matrix over
\(\mathbb{F}\).
\end{definitionbox}

One may therefore ask:

\[
\boxed{
\text{Over which fields is a given matroid representable?}
}
\]

This is a central structural question.

%%%%%%%%%%%%%%%%%%%%%%%%%%%%%%%%%%%%%%%%%%%%%%%%%%%%%%%%%%%%
\subsection{Binary Matroids}
\label{subsec:binary-matroids}
%%%%%%%%%%%%%%%%%%%%%%%%%%%%%%%%%%%%%%%%%%%%%%%%%%%%%%%%%%%%

A matroid representable over

\[
\mathbb{F}_2
\]

is called binary.

Thus a binary matroid can be represented by a matrix whose entries are

\[
0
\]

and

\[
1,
\]

with arithmetic performed modulo \(2\).

Every graphic matroid is binary.

%%%%%%%%%%%%%%%%%%%%%%%%%%%%%%%%%%%%%%%%%%%%%%%%%%%%%%%%%%%%
\subsection{Ternary and Regular Matroids}
\label{subsec:ternary-regular-matroids}
%%%%%%%%%%%%%%%%%%%%%%%%%%%%%%%%%%%%%%%%%%%%%%%%%%%%%%%%%%%%

A matroid representable over

\[
\mathbb{F}_3
\]

is called ternary.

A matroid representable over every field is called regular.

Thus,

\[
\boxed{
\text{regular}
\Longrightarrow
\text{binary and ternary}.
}
\]

Regular matroids occupy an important position between graph theory and
linear algebra.

%%%%%%%%%%%%%%%%%%%%%%%%%%%%%%%%%%%%%%%%%%%%%%%%%%%%%%%%%%%%
\subsection{Uniform Matroids}
\label{subsec:uniform-matroids}
%%%%%%%%%%%%%%%%%%%%%%%%%%%%%%%%%%%%%%%%%%%%%%%%%%%%%%%%%%%%

\begin{definitionbox}{Uniform Matroid}
The uniform matroid

\[
\boxed{
U_{r,n}
}
\]

has ground set of size \(n\), and a subset is independent exactly when
its size is at most \(r\).
\end{definitionbox}

Thus,

\[
\mathcal{I}
=
\{
I\subseteq E:
|I|\leq r
\}.
\]

Every \(r\)-element subset is a basis.

%%%%%%%%%%%%%%%%%%%%%%%%%%%%%%%%%%%%%%%%%%%%%%%%%%%%%%%%%%%%
\subsection{Uniform Matroid Circuits}
\label{subsec:uniform-matroid-circuits}
%%%%%%%%%%%%%%%%%%%%%%%%%%%%%%%%%%%%%%%%%%%%%%%%%%%%%%%%%%%%

In

\[
U_{r,n},
\]

every set of size at most \(r\) is independent.

Therefore the minimal dependent sets have size

\[
r+1.
\]

Hence,

\[
\boxed{
\text{circuits of }U_{r,n}
=
\binom{E}{r+1}.
}
\]

%%%%%%%%%%%%%%%%%%%%%%%%%%%%%%%%%%%%%%%%%%%%%%%%%%%%%%%%%%%%
\subsection{Partition Matroids}
\label{subsec:partition-matroids}
%%%%%%%%%%%%%%%%%%%%%%%%%%%%%%%%%%%%%%%%%%%%%%%%%%%%%%%%%%%%

Suppose the ground set is partitioned as

\[
E
=
E_1\sqcup E_2\sqcup\cdots\sqcup E_m.
\]

Choose nonnegative capacities

\[
r_1,\ldots,r_m.
\]

Declare \(I\subseteq E\) independent when

\[
\boxed{
|I\cap E_j|
\leq
r_j
}
\]

for every \(j\).

This defines a partition matroid.

%%%%%%%%%%%%%%%%%%%%%%%%%%%%%%%%%%%%%%%%%%%%%%%%%%%%%%%%%%%%
\subsection{Worked Example: Partition Matroid}
\label{subsec:worked-partition-matroid}
%%%%%%%%%%%%%%%%%%%%%%%%%%%%%%%%%%%%%%%%%%%%%%%%%%%%%%%%%%%%

\begin{workedexamplebox}{Category Constraints}

Suppose

\[
E=A\sqcup B\sqcup C
\]

and one may select:

\[
\text{at most }2\text{ elements from }A,
\]

\[
\text{at most }1\text{ element from }B,
\]

and

\[
\text{at most }3\text{ elements from }C.
\]

Then the independent sets satisfy

\[
|I\cap A|\leq2,
\]

\[
|I\cap B|\leq1,
\]

and

\[
|I\cap C|\leq3.
\]

\begin{answerbox}
This independence system is a partition matroid of rank

\[
\boxed{
2+1+3=6
}
\]

provided each part contains at least its capacity.
\end{answerbox}

\end{workedexamplebox}

%%%%%%%%%%%%%%%%%%%%%%%%%%%%%%%%%%%%%%%%%%%%%%%%%%%%%%%%%%%%
\subsection{Loops}
\label{subsec:matroid-loops}
%%%%%%%%%%%%%%%%%%%%%%%%%%%%%%%%%%%%%%%%%%%%%%%%%%%%%%%%%%%%

An element

\[
e\in E
\]

is a loop if

\[
\boxed{
\{e\}
}
\]

is dependent.

Thus a loop belongs to no basis.

In a vector matroid, a zero vector corresponds to a loop.

In a graphic matroid, a graph loop corresponds to a matroid loop.

%%%%%%%%%%%%%%%%%%%%%%%%%%%%%%%%%%%%%%%%%%%%%%%%%%%%%%%%%%%%
\subsection{Coloops}
\label{subsec:matroid-coloops}
%%%%%%%%%%%%%%%%%%%%%%%%%%%%%%%%%%%%%%%%%%%%%%%%%%%%%%%%%%%%

An element \(e\) is a coloop if it belongs to every basis.

In a connected graph, a bridge is a coloop of the graphic matroid.

Thus:

\[
\boxed{
\text{graph bridge}
\longleftrightarrow
\text{matroid coloop}.
}
\]

%%%%%%%%%%%%%%%%%%%%%%%%%%%%%%%%%%%%%%%%%%%%%%%%%%%%%%%%%%%%
\subsection{Parallel Elements}
\label{subsec:parallel-elements-matroid}
%%%%%%%%%%%%%%%%%%%%%%%%%%%%%%%%%%%%%%%%%%%%%%%%%%%%%%%%%%%%

Two distinct nonloop elements \(e,f\) are parallel if

\[
\{e,f\}
\]

is a circuit.

In a vector matroid, this means the corresponding nonzero vectors are
scalar multiples.

In a graphic matroid of a multigraph, parallel graph edges produce
parallel matroid elements.

%%%%%%%%%%%%%%%%%%%%%%%%%%%%%%%%%%%%%%%%%%%%%%%%%%%%%%%%%%%%
\subsection{Spanning Sets}
\label{subsec:matroid-spanning}
%%%%%%%%%%%%%%%%%%%%%%%%%%%%%%%%%%%%%%%%%%%%%%%%%%%%%%%%%%%%

A set

\[
S\subseteq E
\]

is spanning if

\[
\boxed{
r(S)=r(E).
}
\]

Thus \(S\) contains enough elements to achieve the full rank of the
matroid.

A basis is simultaneously:

\[
\boxed{
\text{independent}
}
\]

and

\[
\boxed{
\text{spanning}.
}
\]

%%%%%%%%%%%%%%%%%%%%%%%%%%%%%%%%%%%%%%%%%%%%%%%%%%%%%%%%%%%%
\subsection{Closure}
\label{subsec:matroid-closure}
%%%%%%%%%%%%%%%%%%%%%%%%%%%%%%%%%%%%%%%%%%%%%%%%%%%%%%%%%%%%

\begin{definitionbox}{Closure}
For \(A\subseteq E\), the closure of \(A\) is

\[
\boxed{
\operatorname{cl}(A)
=
\{
e\in E:
r(A\cup\{e\})=r(A)
\}.
}
\]
\end{definitionbox}

Thus the closure contains all elements that fail to increase the rank.

This generalizes linear span.

%%%%%%%%%%%%%%%%%%%%%%%%%%%%%%%%%%%%%%%%%%%%%%%%%%%%%%%%%%%%
\subsection{Closure in Vector Matroids}
\label{subsec:closure-vector-matroid}
%%%%%%%%%%%%%%%%%%%%%%%%%%%%%%%%%%%%%%%%%%%%%%%%%%%%%%%%%%%%

For a vector matroid,

\[
e\in\operatorname{cl}(A)
\]

exactly when the vector corresponding to \(e\) lies in the linear span
of the vectors corresponding to \(A\).

Thus,

\[
\boxed{
\text{matroid closure}
\longleftrightarrow
\text{linear span}.
}
\]

%%%%%%%%%%%%%%%%%%%%%%%%%%%%%%%%%%%%%%%%%%%%%%%%%%%%%%%%%%%%
\subsection{Closure in Graphic Matroids}
\label{subsec:closure-graphic-matroid}
%%%%%%%%%%%%%%%%%%%%%%%%%%%%%%%%%%%%%%%%%%%%%%%%%%%%%%%%%%%%

For a graphic matroid,

\[
e=uv
\]

lies in the closure of \(A\) when \(u\) and \(v\) already lie in the
same connected component of the graph \((V,A)\).

Adding \(e\) then creates a cycle and does not increase rank.

%%%%%%%%%%%%%%%%%%%%%%%%%%%%%%%%%%%%%%%%%%%%%%%%%%%%%%%%%%%%
\subsection{Flats}
\label{subsec:matroid-flats}
%%%%%%%%%%%%%%%%%%%%%%%%%%%%%%%%%%%%%%%%%%%%%%%%%%%%%%%%%%%%

\begin{definitionbox}{Flat}
A flat is a set

\[
F\subseteq E
\]

satisfying

\[
\boxed{
\operatorname{cl}(F)=F.
}
\]
\end{definitionbox}

Equivalently, adding any element outside \(F\) increases the rank.

Flats generalize linear subspaces.

%%%%%%%%%%%%%%%%%%%%%%%%%%%%%%%%%%%%%%%%%%%%%%%%%%%%%%%%%%%%
\subsection{Hyperplanes}
\label{subsec:matroid-hyperplanes}
%%%%%%%%%%%%%%%%%%%%%%%%%%%%%%%%%%%%%%%%%%%%%%%%%%%%%%%%%%%%

A hyperplane is a maximal proper flat.

If

\[
r(M)=r,
\]

then every hyperplane has rank

\[
\boxed{
r-1.
}
\]

This generalizes codimension-one subspaces in linear algebra.

%%%%%%%%%%%%%%%%%%%%%%%%%%%%%%%%%%%%%%%%%%%%%%%%%%%%%%%%%%%%
\subsection{Basis Exchange}
\label{subsec:basis-exchange}
%%%%%%%%%%%%%%%%%%%%%%%%%%%%%%%%%%%%%%%%%%%%%%%%%%%%%%%%%%%%

Bases satisfy a powerful exchange rule.

\begin{theorem}[Basis Exchange]
If \(B_1\) and \(B_2\) are bases and

\[
e\in B_1\setminus B_2,
\]

then there exists

\[
f\in B_2\setminus B_1
\]

such that

\[
\boxed{
(B_1-\{e\})\cup\{f\}
}
\]

is also a basis.
\end{theorem}

This is one of the central structural properties of matroids.

%%%%%%%%%%%%%%%%%%%%%%%%%%%%%%%%%%%%%%%%%%%%%%%%%%%%%%%%%%%%
\subsection{Strong Basis Exchange}
\label{subsec:strong-basis-exchange}
%%%%%%%%%%%%%%%%%%%%%%%%%%%%%%%%%%%%%%%%%%%%%%%%%%%%%%%%%%%%

A stronger theorem states that one can choose \(f\) so that both

\[
(B_1-\{e\})\cup\{f\}
\]

and

\[
(B_2-\{f\})\cup\{e\}
\]

are bases.

This reveals a strong symmetry among bases.

%%%%%%%%%%%%%%%%%%%%%%%%%%%%%%%%%%%%%%%%%%%%%%%%%%%%%%%%%%%%
\subsection{Fundamental Circuits}
\label{subsec:fundamental-circuits}
%%%%%%%%%%%%%%%%%%%%%%%%%%%%%%%%%%%%%%%%%%%%%%%%%%%%%%%%%%%%

Let \(B\) be a basis and let

\[
e\notin B.
\]

Then

\[
B\cup\{e\}
\]

is dependent.

It contains a unique circuit, called the fundamental circuit of \(e\)
with respect to \(B\).

Thus,

\[
\boxed{
C(e,B)
\subseteq
B\cup\{e\}.
}
\]

%%%%%%%%%%%%%%%%%%%%%%%%%%%%%%%%%%%%%%%%%%%%%%%%%%%%%%%%%%%%
\subsection{Dual Matroids}
\label{subsec:dual-matroids}
%%%%%%%%%%%%%%%%%%%%%%%%%%%%%%%%%%%%%%%%%%%%%%%%%%%%%%%%%%%%

Matroids possess a natural duality.

\begin{definitionbox}{Dual Matroid}
Let \(M\) be a matroid on ground set \(E\).

The dual matroid

\[
\boxed{
M^*
}
\]

has as bases the complements of the bases of \(M\):

\[
\boxed{
\mathcal{B}(M^*)
=
\{
E\setminus B:
B\in\mathcal{B}(M)
\}.
}
\]
\end{definitionbox}

%%%%%%%%%%%%%%%%%%%%%%%%%%%%%%%%%%%%%%%%%%%%%%%%%%%%%%%%%%%%
\subsection{Rank of the Dual}
\label{subsec:dual-rank}
%%%%%%%%%%%%%%%%%%%%%%%%%%%%%%%%%%%%%%%%%%%%%%%%%%%%%%%%%%%%

If

\[
|E|=n
\]

and

\[
r(M)=r,
\]

then every basis of \(M^*\) has size

\[
n-r.
\]

Therefore,

\[
\boxed{
r(M^*)
=
|E|-r(M).
}
\]

%%%%%%%%%%%%%%%%%%%%%%%%%%%%%%%%%%%%%%%%%%%%%%%%%%%%%%%%%%%%
\subsection{Dual Rank Formula}
\label{subsec:dual-rank-formula}
%%%%%%%%%%%%%%%%%%%%%%%%%%%%%%%%%%%%%%%%%%%%%%%%%%%%%%%%%%%%

For

\[
A\subseteq E,
\]

the dual rank satisfies

\[
\boxed{
r_{M^*}(A)
=
|A|
-r_M(E)
+
r_M(E\setminus A).
}
\]

This formula allows dual properties to be derived from the original
rank function.

%%%%%%%%%%%%%%%%%%%%%%%%%%%%%%%%%%%%%%%%%%%%%%%%%%%%%%%%%%%%
\subsection{Cobases and Cocircuits}
\label{subsec:cobases-cocircuits}
%%%%%%%%%%%%%%%%%%%%%%%%%%%%%%%%%%%%%%%%%%%%%%%%%%%%%%%%%%%%

A basis of the dual matroid is sometimes called a cobasis of \(M\).

A circuit of \(M^*\) is called a cocircuit of \(M\).

Thus,

\[
\boxed{
\text{cocircuit of }M
=
\text{circuit of }M^*.
}
\]

In graphic matroids, cocircuits correspond to minimal edge cuts.

%%%%%%%%%%%%%%%%%%%%%%%%%%%%%%%%%%%%%%%%%%%%%%%%%%%%%%%%%%%%
\subsection{Graphic Duality}
\label{subsec:graphic-duality}
%%%%%%%%%%%%%%%%%%%%%%%%%%%%%%%%%%%%%%%%%%%%%%%%%%%%%%%%%%%%

For a planar graph \(G\), the planar dual graph \(G^*\) has:

\begin{itemize}
    \item one vertex for each face of \(G\);
    \item one dual edge crossing each edge of \(G\).
\end{itemize}

Under appropriate conditions,

\[
\boxed{
M(G)^*
\cong
M(G^*).
}
\]

Thus planar graph duality is a special case of matroid duality.

%%%%%%%%%%%%%%%%%%%%%%%%%%%%%%%%%%%%%%%%%%%%%%%%%%%%%%%%%%%%
\subsection{Deletion}
\label{subsec:matroid-deletion}
%%%%%%%%%%%%%%%%%%%%%%%%%%%%%%%%%%%%%%%%%%%%%%%%%%%%%%%%%%%%

Let

\[
e\in E.
\]

The deletion of \(e\) is denoted

\[
\boxed{
M\setminus e.
}
\]

Its ground set is

\[
E\setminus\{e\}.
\]

Its independent sets are simply the independent sets of \(M\) that do
not contain \(e\).

%%%%%%%%%%%%%%%%%%%%%%%%%%%%%%%%%%%%%%%%%%%%%%%%%%%%%%%%%%%%
\subsection{Contraction}
\label{subsec:matroid-contraction}
%%%%%%%%%%%%%%%%%%%%%%%%%%%%%%%%%%%%%%%%%%%%%%%%%%%%%%%%%%%%

Contraction is the dual operation to deletion.

The contraction of \(e\) is denoted

\[
\boxed{
M/e.
}
\]

For a nonloop element \(e\), the rank function satisfies

\[
\boxed{
r_{M/e}(A)
=
r_M(A\cup\{e\})-r_M(\{e\})
}
\]

for

\[
A\subseteq E\setminus\{e\}.
\]

For a nonloop,

\[
r_M(\{e\})=1.
\]

%%%%%%%%%%%%%%%%%%%%%%%%%%%%%%%%%%%%%%%%%%%%%%%%%%%%%%%%%%%%
\subsection{Deletion--Contraction Duality}
\label{subsec:deletion-contraction-duality}
%%%%%%%%%%%%%%%%%%%%%%%%%%%%%%%%%%%%%%%%%%%%%%%%%%%%%%%%%%%%

Deletion and contraction are exchanged by duality:

\[
\boxed{
(M\setminus e)^*
=
M^*/e
}
\]

and

\[
\boxed{
(M/e)^*
=
M^*\setminus e.
}
\]

This is one reason deletion--contraction recurrences appear naturally
throughout graph and matroid theory.

%%%%%%%%%%%%%%%%%%%%%%%%%%%%%%%%%%%%%%%%%%%%%%%%%%%%%%%%%%%%
\subsection{Graphic Deletion and Contraction}
\label{subsec:graphic-deletion-contraction}
%%%%%%%%%%%%%%%%%%%%%%%%%%%%%%%%%%%%%%%%%%%%%%%%%%%%%%%%%%%%

For a graphic matroid:

\[
M(G)\setminus e
\]

corresponds to deleting the graph edge \(e\).

Similarly,

\[
M(G)/e
\]

corresponds to contracting the graph edge \(e\), with the usual care
for loops and parallel edges.

Thus abstract matroid operations generalize familiar graph operations.

%%%%%%%%%%%%%%%%%%%%%%%%%%%%%%%%%%%%%%%%%%%%%%%%%%%%%%%%%%%%
\subsection{Minors}
\label{subsec:matroid-minors}
%%%%%%%%%%%%%%%%%%%%%%%%%%%%%%%%%%%%%%%%%%%%%%%%%%%%%%%%%%%%

A matroid obtained through a sequence of deletions and contractions is
called a minor.

Thus,

\[
\boxed{
\text{matroid minor}
=
\text{repeated deletion and contraction}.
}
\]

Minor-closed classes of matroids play a role analogous to minor-closed
classes of graphs.

%%%%%%%%%%%%%%%%%%%%%%%%%%%%%%%%%%%%%%%%%%%%%%%%%%%%%%%%%%%%
\subsection{Excluded Minors}
\label{subsec:excluded-minors-matroid}
%%%%%%%%%%%%%%%%%%%%%%%%%%%%%%%%%%%%%%%%%%%%%%%%%%%%%%%%%%%%

A class of matroids may sometimes be characterized by forbidden minors.

An excluded minor for a class is a matroid not in the class, all of
whose proper minors belong to the class.

This generalizes forbidden-minor characterizations in graph theory.

%%%%%%%%%%%%%%%%%%%%%%%%%%%%%%%%%%%%%%%%%%%%%%%%%%%%%%%%%%%%
\subsection{Representability and Minors}
\label{subsec:representability-minors}
%%%%%%%%%%%%%%%%%%%%%%%%%%%%%%%%%%%%%%%%%%%%%%%%%%%%%%%%%%%%

If a matroid is representable over a field \(\mathbb{F}\), then every
minor is also representable over \(\mathbb{F}\).

Therefore representability over a fixed field is a minor-closed
property.

This makes excluded-minor characterizations possible.

%%%%%%%%%%%%%%%%%%%%%%%%%%%%%%%%%%%%%%%%%%%%%%%%%%%%%%%%%%%%
\subsection{The Fano Matroid}
\label{subsec:fano-matroid}
%%%%%%%%%%%%%%%%%%%%%%%%%%%%%%%%%%%%%%%%%%%%%%%%%%%%%%%%%%%%

The Fano plane produces an important rank-\(3\) matroid on seven
elements.

Its elements are the seven points of the Fano plane.

A triple is dependent exactly when its three points lie on a Fano line.

This matroid is denoted

\[
\boxed{
F_7.
}
\]

%%%%%%%%%%%%%%%%%%%%%%%%%%%%%%%%%%%%%%%%%%%%%%%%%%%%%%%%%%%%
\subsection{Representation of the Fano Matroid}
\label{subsec:fano-matroid-representation}
%%%%%%%%%%%%%%%%%%%%%%%%%%%%%%%%%%%%%%%%%%%%%%%%%%%%%%%%%%%%

The Fano matroid can be represented over

\[
\mathbb{F}_2
\]

using the seven nonzero vectors of

\[
\mathbb{F}_2^3.
\]

For example, use columns

\[
\begin{pmatrix}
1\\0\\0
\end{pmatrix},
\begin{pmatrix}
0\\1\\0
\end{pmatrix},
\begin{pmatrix}
0\\0\\1
\end{pmatrix},
\begin{pmatrix}
1\\1\\0
\end{pmatrix},
\begin{pmatrix}
1\\0\\1
\end{pmatrix},
\begin{pmatrix}
0\\1\\1
\end{pmatrix},
\begin{pmatrix}
1\\1\\1
\end{pmatrix}.
\]

Thus,

\[
\boxed{
F_7
\text{ is binary}.
}
\]

%%%%%%%%%%%%%%%%%%%%%%%%%%%%%%%%%%%%%%%%%%%%%%%%%%%%%%%%%%%%
\subsection{Projective Geometry Matroids}
\label{subsec:projective-geometry-matroids}
%%%%%%%%%%%%%%%%%%%%%%%%%%%%%%%%%%%%%%%%%%%%%%%%%%%%%%%%%%%%

The points of projective geometry

\[
\operatorname{PG}(r-1,q)
\]

define a rank-\(r\) matroid.

Independence is inherited from linear independence of representative
vectors in

\[
\mathbb{F}_q^r.
\]

This produces a natural bridge among:

\[
\boxed{
\text{finite geometry}
+
\text{linear algebra}
+
\text{design theory}
+
\text{matroids}.
}
\]

%%%%%%%%%%%%%%%%%%%%%%%%%%%%%%%%%%%%%%%%%%%%%%%%%%%%%%%%%%%%
\subsection{Affine Geometry Matroids}
\label{subsec:affine-geometry-matroids}
%%%%%%%%%%%%%%%%%%%%%%%%%%%%%%%%%%%%%%%%%%%%%%%%%%%%%%%%%%%%

Finite affine geometries similarly produce representable matroids.

Affine dependence may be converted into linear dependence by using
homogeneous coordinates.

Thus many finite geometric incidence structures possess a natural
matroid interpretation.

%%%%%%%%%%%%%%%%%%%%%%%%%%%%%%%%%%%%%%%%%%%%%%%%%%%%%%%%%%%%
\subsection{Greedy Algorithms}
\label{subsec:matroid-greedy}
%%%%%%%%%%%%%%%%%%%%%%%%%%%%%%%%%%%%%%%%%%%%%%%%%%%%%%%%%%%%

One of the deepest elementary properties of matroids is their
relationship with greedy optimization.

Suppose every element

\[
e\in E
\]

has a weight

\[
w(e).
\]

We want an independent set of maximum total weight.

%%%%%%%%%%%%%%%%%%%%%%%%%%%%%%%%%%%%%%%%%%%%%%%%%%%%%%%%%%%%
\subsection{Matroid Greedy Algorithm}
\label{subsec:matroid-greedy-algorithm}
%%%%%%%%%%%%%%%%%%%%%%%%%%%%%%%%%%%%%%%%%%%%%%%%%%%%%%%%%%%%

Order the elements so that

\[
w(e_1)\geq w(e_2)\geq\cdots\geq w(e_n).
\]

Start with

\[
I=\varnothing.
\]

For \(j=1,\ldots,n\):

\begin{itemize}
    \item if
    \[
    I\cup\{e_j\}
    \]
    is independent, add \(e_j\);
    \item otherwise skip \(e_j\).
\end{itemize}

The final set is a maximum-weight basis when the weights are
nonnegative, and more generally yields the appropriate maximum-weight
independent solution under the usual formulation.

%%%%%%%%%%%%%%%%%%%%%%%%%%%%%%%%%%%%%%%%%%%%%%%%%%%%%%%%%%%%
\subsection{Greedy Characterization of Matroids}
\label{subsec:greedy-characterization-matroids}
%%%%%%%%%%%%%%%%%%%%%%%%%%%%%%%%%%%%%%%%%%%%%%%%%%%%%%%%%%%%

\begin{theorem}[Greedy Characterization]
A finite hereditary independence system is a matroid if and only if the
greedy algorithm produces a maximum-weight basis for every assignment
of weights.
\end{theorem}

Thus matroids are exactly the independence systems for which the
natural greedy algorithm always works.

This gives a powerful algorithmic interpretation of the exchange axiom.

%%%%%%%%%%%%%%%%%%%%%%%%%%%%%%%%%%%%%%%%%%%%%%%%%%%%%%%%%%%%
\subsection{Worked Example: Greedy in a Graphic Matroid}
\label{subsec:worked-greedy-graphic}
%%%%%%%%%%%%%%%%%%%%%%%%%%%%%%%%%%%%%%%%%%%%%%%%%%%%%%%%%%%%

\begin{workedexamplebox}{Minimum Spanning Tree Connection}

Let \(G\) be a connected weighted graph.

The independent sets of the graphic matroid are forests.

To obtain a minimum-weight spanning tree, process edges from lowest to
highest weight and add an edge whenever doing so does not create a
cycle.

This is Kruskal's algorithm.

The matroid exchange property guarantees optimality.

\begin{answerbox}
\[
\boxed{
\text{Kruskal's algorithm is a matroid greedy algorithm.}
}
\]
\end{answerbox}

\end{workedexamplebox}

%%%%%%%%%%%%%%%%%%%%%%%%%%%%%%%%%%%%%%%%%%%%%%%%%%%%%%%%%%%%
\subsection{Maximum-Weight Bases}
\label{subsec:maximum-weight-bases}
%%%%%%%%%%%%%%%%%%%%%%%%%%%%%%%%%%%%%%%%%%%%%%%%%%%%%%%%%%%%

For a weight function

\[
w:E\to\mathbb{R},
\]

the weight of a set is

\[
\boxed{
w(A)
=
\sum_{e\in A}w(e).
}
\]

The maximum-weight basis problem is

\[
\boxed{
\max
\{
w(B):
B\text{ is a basis}
\}.
}
\]

For matroids, this optimization problem is solved greedily.

%%%%%%%%%%%%%%%%%%%%%%%%%%%%%%%%%%%%%%%%%%%%%%%%%%%%%%%%%%%%
\subsection{Matroid Polytopes Preview}
\label{subsec:matroid-polytopes-preview}
%%%%%%%%%%%%%%%%%%%%%%%%%%%%%%%%%%%%%%%%%%%%%%%%%%%%%%%%%%%%

Associate to every basis \(B\) its incidence vector

\[
\mathbf{1}_B\in\mathbb{R}^{E}.
\]

The convex hull

\[
\boxed{
P(M)
=
\operatorname{conv}
\{
\mathbf{1}_B:
B\text{ basis}
\}
}
\]

is called the base polytope of the matroid.

Matroid polytopes connect discrete independence with convex
optimization.

%%%%%%%%%%%%%%%%%%%%%%%%%%%%%%%%%%%%%%%%%%%%%%%%%%%%%%%%%%%%
\subsection{Independent-Set Polytope}
\label{subsec:independent-set-polytope}
%%%%%%%%%%%%%%%%%%%%%%%%%%%%%%%%%%%%%%%%%%%%%%%%%%%%%%%%%%%%

The convex hull of the incidence vectors of all independent sets is the
independence polytope.

It can be described using rank inequalities:

\[
\boxed{
\sum_{e\in A}x_e
\leq
r(A)
}
\]

for every

\[
A\subseteq E,
\]

together with

\[
x_e\geq0.
\]

This is a fundamental bridge between submodularity and optimization.

%%%%%%%%%%%%%%%%%%%%%%%%%%%%%%%%%%%%%%%%%%%%%%%%%%%%%%%%%%%%
\subsection{Matroid Intersection}
\label{subsec:matroid-intersection}
%%%%%%%%%%%%%%%%%%%%%%%%%%%%%%%%%%%%%%%%%%%%%%%%%%%%%%%%%%%%

Suppose

\[
M_1=(E,\mathcal{I}_1)
\]

and

\[
M_2=(E,\mathcal{I}_2)
\]

are matroids on the same ground set.

The matroid intersection problem seeks

\[
\boxed{
\max
\{
|I|:
I\in\mathcal{I}_1\cap\mathcal{I}_2
\}.
}
\]

Unlike optimization over one matroid, simple greedy selection is not
generally sufficient.

%%%%%%%%%%%%%%%%%%%%%%%%%%%%%%%%%%%%%%%%%%%%%%%%%%%%%%%%%%%%
\subsection{Bipartite Matching as Matroid Intersection}
\label{subsec:matching-matroid-intersection}
%%%%%%%%%%%%%%%%%%%%%%%%%%%%%%%%%%%%%%%%%%%%%%%%%%%%%%%%%%%%

A bipartite matching can be viewed using two partition matroids.

Let the ground set consist of graph edges.

One matroid permits at most one chosen edge incident with each vertex
in the left part.

The second permits at most one chosen edge incident with each vertex in
the right part.

A set independent in both matroids is exactly a matching.

Thus,

\[
\boxed{
\text{bipartite matching}
=
\text{intersection of two partition matroids}.
}
\]

%%%%%%%%%%%%%%%%%%%%%%%%%%%%%%%%%%%%%%%%%%%%%%%%%%%%%%%%%%%%
\subsection{Matroid Intersection Min--Max Theorem}
\label{subsec:matroid-intersection-theorem}
%%%%%%%%%%%%%%%%%%%%%%%%%%%%%%%%%%%%%%%%%%%%%%%%%%%%%%%%%%%%

A fundamental theorem gives a min--max characterization of maximum
common independent sets.

For matroids \(M_1\) and \(M_2\) on \(E\),

\[
\boxed{
\max
\{
|I|:
I\in\mathcal{I}_1\cap\mathcal{I}_2
\}
=
\min_{A\subseteq E}
\left(
r_1(A)+r_2(E\setminus A)
\right).
}
\]

This generalizes several classical matching and covering theorems.

%%%%%%%%%%%%%%%%%%%%%%%%%%%%%%%%%%%%%%%%%%%%%%%%%%%%%%%%%%%%
\subsection{Matroid Union}
\label{subsec:matroid-union}
%%%%%%%%%%%%%%%%%%%%%%%%%%%%%%%%%%%%%%%%%%%%%%%%%%%%%%%%%%%%

Given matroids

\[
M_1,\ldots,M_k
\]

on the same ground set \(E\), matroid union studies sets expressible as

\[
I_1\cup\cdots\cup I_k,
\]

where each

\[
I_j
\]

is independent in \(M_j\).

This generalizes partitioning structures into several independent
pieces.

%%%%%%%%%%%%%%%%%%%%%%%%%%%%%%%%%%%%%%%%%%%%%%%%%%%%%%%%%%%%
\subsection{Arboricity Connection}
\label{subsec:matroid-union-arboricity}
%%%%%%%%%%%%%%%%%%%%%%%%%%%%%%%%%%%%%%%%%%%%%%%%%%%%%%%%%%%%

For a graphic matroid, matroid union asks whether a graph's edges can be
partitioned into forests.

This leads naturally to graph arboricity.

The arboricity of a graph is the minimum number of forests whose union
contains all graph edges.

Thus matroid union generalizes a classical graph decomposition problem.

%%%%%%%%%%%%%%%%%%%%%%%%%%%%%%%%%%%%%%%%%%%%%%%%%%%%%%%%%%%%
\subsection{Transversal Matroids}
\label{subsec:transversal-matroids}
%%%%%%%%%%%%%%%%%%%%%%%%%%%%%%%%%%%%%%%%%%%%%%%%%%%%%%%%%%%%

Let

\[
\mathcal{A}
=
(A_1,\ldots,A_m)
\]

be a family of subsets of \(E\).

A partial transversal is a set of distinct representatives chosen from
distinct members of \(\mathcal{A}\).

The subsets of \(E\) that can occur as partial transversals form the
independent sets of a matroid.

This is called a transversal matroid.

%%%%%%%%%%%%%%%%%%%%%%%%%%%%%%%%%%%%%%%%%%%%%%%%%%%%%%%%%%%%
\subsection{Hall's Theorem Connection}
\label{subsec:transversal-hall}
%%%%%%%%%%%%%%%%%%%%%%%%%%%%%%%%%%%%%%%%%%%%%%%%%%%%%%%%%%%%

Transversal matroids are closely related to bipartite matching.

Construct a bipartite graph with:

\begin{itemize}
    \item one side representing the sets \(A_i\);
    \item the other side representing elements of \(E\);
    \item an edge when an element belongs to a set.
\end{itemize}

Partial transversals correspond to matchings.

Thus Hall's Marriage Theorem naturally enters matroid theory.

%%%%%%%%%%%%%%%%%%%%%%%%%%%%%%%%%%%%%%%%%%%%%%%%%%%%%%%%%%%%
\subsection{Gammoids Preview}
\label{subsec:gammoids-preview}
%%%%%%%%%%%%%%%%%%%%%%%%%%%%%%%%%%%%%%%%%%%%%%%%%%%%%%%%%%%%

A gammoid is a matroid derived from vertex-disjoint directed paths in a
digraph.

A set is independent when its elements can be linked by disjoint paths
to a specified target set.

Gammoids extend the connection between matroids and network flow.

%%%%%%%%%%%%%%%%%%%%%%%%%%%%%%%%%%%%%%%%%%%%%%%%%%%%%%%%%%%%
\subsection{Direct Sums}
\label{subsec:matroid-direct-sums}
%%%%%%%%%%%%%%%%%%%%%%%%%%%%%%%%%%%%%%%%%%%%%%%%%%%%%%%%%%%%

Let

\[
M_1=(E_1,\mathcal{I}_1)
\]

and

\[
M_2=(E_2,\mathcal{I}_2)
\]

with

\[
E_1\cap E_2=\varnothing.
\]

Their direct sum

\[
\boxed{
M_1\oplus M_2
}
\]

has ground set

\[
E_1\cup E_2
\]

and independent sets

\[
I_1\cup I_2
\]

with

\[
I_j\in\mathcal{I}_j.
\]

%%%%%%%%%%%%%%%%%%%%%%%%%%%%%%%%%%%%%%%%%%%%%%%%%%%%%%%%%%%%
\subsection{Connected Matroids}
\label{subsec:connected-matroids}
%%%%%%%%%%%%%%%%%%%%%%%%%%%%%%%%%%%%%%%%%%%%%%%%%%%%%%%%%%%%

A matroid is connected if it cannot be written as a nontrivial direct
sum.

For graphic matroids, matroid connectedness is related to the cycle
structure of the graph rather than merely ordinary graph connectedness.

Thus matroid connectivity captures dependence connectivity.

%%%%%%%%%%%%%%%%%%%%%%%%%%%%%%%%%%%%%%%%%%%%%%%%%%%%%%%%%%%%
\subsection{Restriction}
\label{subsec:matroid-restriction}
%%%%%%%%%%%%%%%%%%%%%%%%%%%%%%%%%%%%%%%%%%%%%%%%%%%%%%%%%%%%

For

\[
A\subseteq E,
\]

the restriction of \(M\) to \(A\), written

\[
\boxed{
M|A,
}
\]

is the matroid whose independent sets are the independent subsets of
\(A\).

Restriction is equivalent to deleting every element outside \(A\).

%%%%%%%%%%%%%%%%%%%%%%%%%%%%%%%%%%%%%%%%%%%%%%%%%%%%%%%%%%%%
\subsection{Truncation}
\label{subsec:matroid-truncation}
%%%%%%%%%%%%%%%%%%%%%%%%%%%%%%%%%%%%%%%%%%%%%%%%%%%%%%%%%%%%

Suppose \(M\) has rank \(r\).

Its truncation lowers the maximum allowed independent-set size by one.

Thus the truncated matroid has independent sets

\[
\boxed{
\{
I\in\mathcal{I}(M):
|I|\leq r-1
\}.
}
\]

More generally, one may truncate to any smaller rank.

%%%%%%%%%%%%%%%%%%%%%%%%%%%%%%%%%%%%%%%%%%%%%%%%%%%%%%%%%%%%
\subsection{Tutte Polynomial of a Matroid}
\label{subsec:matroid-tutte-polynomial}
%%%%%%%%%%%%%%%%%%%%%%%%%%%%%%%%%%%%%%%%%%%%%%%%%%%%%%%%%%%%

The Tutte polynomial extends naturally from graphs to matroids.

For a matroid \(M=(E,r)\),

\[
\boxed{
T_M(x,y)
=
\sum_{A\subseteq E}
(x-1)^{r(E)-r(A)}
(y-1)^{|A|-r(A)}.
}
\]

This polynomial encodes substantial structural information.

%%%%%%%%%%%%%%%%%%%%%%%%%%%%%%%%%%%%%%%%%%%%%%%%%%%%%%%%%%%%
\subsection{Rank and Nullity in the Tutte Polynomial}
\label{subsec:rank-nullity-tutte}
%%%%%%%%%%%%%%%%%%%%%%%%%%%%%%%%%%%%%%%%%%%%%%%%%%%%%%%%%%%%

The exponent

\[
r(E)-r(A)
\]

measures how far \(A\) is from spanning.

The exponent

\[
|A|-r(A)
\]

is the nullity of \(A\), measuring how much dependence it contains.

Thus the Tutte polynomial simultaneously records spanning deficiency
and dependence.

%%%%%%%%%%%%%%%%%%%%%%%%%%%%%%%%%%%%%%%%%%%%%%%%%%%%%%%%%%%%
\subsection{Deletion--Contraction for the Tutte Polynomial}
\label{subsec:tutte-deletion-contraction-matroid}
%%%%%%%%%%%%%%%%%%%%%%%%%%%%%%%%%%%%%%%%%%%%%%%%%%%%%%%%%%%%

If \(e\) is neither a loop nor a coloop, then

\[
\boxed{
T_M(x,y)
=
T_{M\setminus e}(x,y)
+
T_{M/e}(x,y).
}
\]

If \(e\) is a coloop,

\[
\boxed{
T_M(x,y)
=
xT_{M/e}(x,y).
}
\]

If \(e\) is a loop,

\[
\boxed{
T_M(x,y)
=
yT_{M\setminus e}(x,y).
}
\]

%%%%%%%%%%%%%%%%%%%%%%%%%%%%%%%%%%%%%%%%%%%%%%%%%%%%%%%%%%%%
\subsection{Tutte Polynomial and Bases}
\label{subsec:tutte-bases}
%%%%%%%%%%%%%%%%%%%%%%%%%%%%%%%%%%%%%%%%%%%%%%%%%%%%%%%%%%%%

A particularly important evaluation is

\[
\boxed{
T_M(1,1)
=
\text{number of bases of }M.
}
\]

For a connected graphic matroid,

\[
T_{M(G)}(1,1)
\]

therefore equals the number of spanning trees of \(G\).

%%%%%%%%%%%%%%%%%%%%%%%%%%%%%%%%%%%%%%%%%%%%%%%%%%%%%%%%%%%%
\subsection{Characteristic Polynomial Preview}
\label{subsec:matroid-characteristic-polynomial}
%%%%%%%%%%%%%%%%%%%%%%%%%%%%%%%%%%%%%%%%%%%%%%%%%%%%%%%%%%%%

A matroid also has a characteristic polynomial.

For loopless \(M\), one form is

\[
\boxed{
\chi_M(\lambda)
=
\sum_{A\subseteq E}
(-1)^{|A|}
\lambda^{r(E)-r(A)}.
}
\]

This generalizes graph chromatic phenomena and is closely related to
the lattice of flats.

%%%%%%%%%%%%%%%%%%%%%%%%%%%%%%%%%%%%%%%%%%%%%%%%%%%%%%%%%%%%
\subsection{Lattice of Flats}
\label{subsec:lattice-of-flats}
%%%%%%%%%%%%%%%%%%%%%%%%%%%%%%%%%%%%%%%%%%%%%%%%%%%%%%%%%%%%

The flats of a matroid, ordered by inclusion, form a lattice.

This is called the lattice of flats.

For representable matroids, this resembles the lattice of linear
subspaces.

For graphic matroids, flats correspond to edge sets closed under
connectivity.

Thus matroid theory naturally connects with lattice theory.

%%%%%%%%%%%%%%%%%%%%%%%%%%%%%%%%%%%%%%%%%%%%%%%%%%%%%%%%%%%%
\subsection{Geometric Lattices}
\label{subsec:geometric-lattices}
%%%%%%%%%%%%%%%%%%%%%%%%%%%%%%%%%%%%%%%%%%%%%%%%%%%%%%%%%%%%

The lattice of flats of a simple matroid is a geometric lattice.

Conversely, finite geometric lattices correspond to simple matroids.

This provides a deep equivalence between:

\[
\boxed{
\text{matroids}
\quad\text{and}\quad
\text{certain lattices}.
}
\]

%%%%%%%%%%%%%%%%%%%%%%%%%%%%%%%%%%%%%%%%%%%%%%%%%%%%%%%%%%%%
\subsection{Simple Matroids}
\label{subsec:simple-matroids}
%%%%%%%%%%%%%%%%%%%%%%%%%%%%%%%%%%%%%%%%%%%%%%%%%%%%%%%%%%%%

A matroid is simple if it contains:

\begin{itemize}
    \item no loops;
    \item no parallel elements.
\end{itemize}

Thus every circuit has size at least \(3\).

Simple matroids are analogous to simple graphs.

%%%%%%%%%%%%%%%%%%%%%%%%%%%%%%%%%%%%%%%%%%%%%%%%%%%%%%%%%%%%
\subsection{Cosimple Matroids}
\label{subsec:cosimple-matroids}
%%%%%%%%%%%%%%%%%%%%%%%%%%%%%%%%%%%%%%%%%%%%%%%%%%%%%%%%%%%%

A matroid is cosimple if its dual is simple.

Equivalently, it has:

\begin{itemize}
    \item no coloops;
    \item no two-element cocircuits.
\end{itemize}

Thus simplicity and cosimplicity are dual notions.

%%%%%%%%%%%%%%%%%%%%%%%%%%%%%%%%%%%%%%%%%%%%%%%%%%%%%%%%%%%%
\subsection{Matroid Isomorphism}
\label{subsec:matroid-isomorphism}
%%%%%%%%%%%%%%%%%%%%%%%%%%%%%%%%%%%%%%%%%%%%%%%%%%%%%%%%%%%%

Matroids

\[
M_1=(E_1,\mathcal{I}_1)
\]

and

\[
M_2=(E_2,\mathcal{I}_2)
\]

are isomorphic if there exists a bijection

\[
f:E_1\to E_2
\]

such that

\[
\boxed{
I\in\mathcal{I}_1
\iff
f(I)\in\mathcal{I}_2.
}
\]

Thus matroid isomorphisms preserve independence structure.

%%%%%%%%%%%%%%%%%%%%%%%%%%%%%%%%%%%%%%%%%%%%%%%%%%%%%%%%%%%%
\subsection{Automorphisms of Matroids}
\label{subsec:matroid-automorphisms}
%%%%%%%%%%%%%%%%%%%%%%%%%%%%%%%%%%%%%%%%%%%%%%%%%%%%%%%%%%%%

An automorphism is a matroid isomorphism from a matroid to itself.

The automorphisms form a group.

For representable or geometric matroids, these symmetries may arise
from:

\begin{itemize}
    \item graph automorphisms;
    \item linear transformations;
    \item projective transformations;
    \item design automorphisms.
\end{itemize}

%%%%%%%%%%%%%%%%%%%%%%%%%%%%%%%%%%%%%%%%%%%%%%%%%%%%%%%%%%%%
\subsection{Matroids and Coding Theory}
\label{subsec:matroids-coding-theory}
%%%%%%%%%%%%%%%%%%%%%%%%%%%%%%%%%%%%%%%%%%%%%%%%%%%%%%%%%%%%

A linear code defined by a matrix determines a representable matroid
from the matrix columns.

Linear dependence among columns records structural properties of the
code.

Circuits correspond to minimal dependent column sets.

This links matroid invariants with:

\begin{itemize}
    \item parity-check matrices;
    \item minimum distance;
    \item code supports;
    \item finite geometry.
\end{itemize}

%%%%%%%%%%%%%%%%%%%%%%%%%%%%%%%%%%%%%%%%%%%%%%%%%%%%%%%%%%%%
\subsection{Matroids and Design Theory}
\label{subsec:matroids-design-theory}
%%%%%%%%%%%%%%%%%%%%%%%%%%%%%%%%%%%%%%%%%%%%%%%%%%%%%%%%%%%%

Finite projective and affine geometries support both:

\[
\text{combinatorial designs}
\]

and

\[
\text{representable matroids}.
\]

For example, the Fano plane gives:

\[
\boxed{
\text{a }2\text{-}(7,3,1)\text{ design}
}
\]

and simultaneously

\[
\boxed{
\text{the rank-3 binary Fano matroid }F_7.
}
\]

Thus the same incidence structure can be interpreted through different
mathematical frameworks.

%%%%%%%%%%%%%%%%%%%%%%%%%%%%%%%%%%%%%%%%%%%%%%%%%%%%%%%%%%%%
\subsection{Matroids and Optimization}
\label{subsec:matroids-optimization}
%%%%%%%%%%%%%%%%%%%%%%%%%%%%%%%%%%%%%%%%%%%%%%%%%%%%%%%%%%%%

Matroids provide one of the clearest settings in which combinatorial
structure guarantees efficient optimization.

Important problems include:

\[
\boxed{
\text{maximum-weight basis},
}
\]

\[
\boxed{
\text{matroid intersection},
}
\]

\[
\boxed{
\text{matroid union},
}
\]

and

\[
\boxed{
\text{submodular optimization}.
}
\]

These ideas are developed more algorithmically in optimization and
operations research.

%%%%%%%%%%%%%%%%%%%%%%%%%%%%%%%%%%%%%%%%%%%%%%%%%%%%%%%%%%%%
\subsection{Why the Exchange Axiom Matters}
\label{subsec:why-exchange-matters}
%%%%%%%%%%%%%%%%%%%%%%%%%%%%%%%%%%%%%%%%%%%%%%%%%%%%%%%%%%%%

Without the exchange axiom, maximal independent sets can have different
sizes.

Greedy choices can then become trapped in a small maximal set.

The exchange property prevents this.

It guarantees that local independent choices can be exchanged and
extended toward globally optimal bases.

Thus,

\[
\boxed{
\text{exchange}
\Longrightarrow
\text{uniform basis size}
+
\text{greedy optimality}.
}
\]

%%%%%%%%%%%%%%%%%%%%%%%%%%%%%%%%%%%%%%%%%%%%%%%%%%%%%%%%%%%%
\subsection{Independence Systems That Are Not Matroids}
\label{subsec:nonmatroid-independence-systems}
%%%%%%%%%%%%%%%%%%%%%%%%%%%%%%%%%%%%%%%%%%%%%%%%%%%%%%%%%%%%

Not every hereditary family is a matroid.

Consider

\[
E=\{1,2,3\}
\]

with independent sets

\[
\varnothing,
\{1\},
\{2\},
\{3\},
\{1,2\}.
\]

This family is hereditary.

However,

\[
I=\{3\}
\]

and

\[
J=\{1,2\}
\]

satisfy

\[
|I|<|J|,
\]

but neither

\[
\{1,3\}
\]

nor

\[
\{2,3\}
\]

is independent.

Thus exchange fails.

Therefore this independence system is not a matroid.

%%%%%%%%%%%%%%%%%%%%%%%%%%%%%%%%%%%%%%%%%%%%%%%%%%%%%%%%%%%%
\subsection{Worked Example: Testing a Matroid}
\label{subsec:worked-testing-matroid}
%%%%%%%%%%%%%%%%%%%%%%%%%%%%%%%%%%%%%%%%%%%%%%%%%%%%%%%%%%%%

\begin{workedexamplebox}{Checking the Exchange Axiom}

Let

\[
E=\{a,b,c\}
\]

and let

\[
\mathcal{I}
=
\{
\varnothing,
\{a\},
\{b\},
\{c\},
\{a,b\},
\{a,c\},
\{b,c\}
\}.
\]

Every subset of an independent set is independent.

Now take any independent sets \(I,J\) with

\[
|I|<|J|.
\]

Because every two-element subset is independent, one can always add an
element of \(J\setminus I\) when \(I\) has size \(0\) or \(1\).

Thus the exchange axiom holds.

\begin{answerbox}
\[
\boxed{
M\cong U_{2,3}.
}
\]
\end{answerbox}

\end{workedexamplebox}

%%%%%%%%%%%%%%%%%%%%%%%%%%%%%%%%%%%%%%%%%%%%%%%%%%%%%%%%%%%%
\subsection{Common Errors}
\label{subsec:matroid-common-errors}
%%%%%%%%%%%%%%%%%%%%%%%%%%%%%%%%%%%%%%%%%%%%%%%%%%%%%%%%%%%%

\subsubsection{Confusing Maximal and Maximum}

A maximal independent set cannot be enlarged while remaining
independent.

A maximum independent set has largest possible cardinality.

In a matroid, every maximal independent set is maximum.

In a general independence system, this need not hold.

%%%%%%%%%%%%%%%%%%%%%%%%%%%%%%%%%%%%%%%%%%%%%%%%%%%%%%%%%%%%
\subsubsection{Assuming Every Hereditary System Is a Matroid}

Hereditary independence is necessary but not sufficient.

The exchange axiom must also hold.

%%%%%%%%%%%%%%%%%%%%%%%%%%%%%%%%%%%%%%%%%%%%%%%%%%%%%%%%%%%%
\subsubsection{Confusing Circuits with Arbitrary Dependent Sets}

A circuit is minimally dependent.

A larger dependent set containing a circuit is not itself necessarily a
circuit.

%%%%%%%%%%%%%%%%%%%%%%%%%%%%%%%%%%%%%%%%%%%%%%%%%%%%%%%%%%%%
\subsubsection{Confusing Graph Circuits with All Graph Cycles in General Matroids}

For graphic matroids, circuits are graph cycles.

For arbitrary matroids, circuits are abstract minimal dependent sets.

%%%%%%%%%%%%%%%%%%%%%%%%%%%%%%%%%%%%%%%%%%%%%%%%%%%%%%%%%%%%
\subsubsection{Confusing Rank with Cardinality}

For a dependent set \(A\),

\[
r(A)<|A|.
\]

Rank measures the maximum size of an independent subset, not the number
of elements present.

%%%%%%%%%%%%%%%%%%%%%%%%%%%%%%%%%%%%%%%%%%%%%%%%%%%%%%%%%%%%
\subsubsection{Confusing Basis and Spanning Set}

Every basis is spanning.

A spanning set may contain dependent elements and therefore may be
larger than a basis.

%%%%%%%%%%%%%%%%%%%%%%%%%%%%%%%%%%%%%%%%%%%%%%%%%%%%%%%%%%%%
\subsubsection{Confusing Closure with Set-Theoretic Closure}

Matroid closure is defined through rank:

\[
e\in\operatorname{cl}(A)
\iff
r(A\cup\{e\})=r(A).
\]

It is an abstract dependence operator, not a topological closure.

%%%%%%%%%%%%%%%%%%%%%%%%%%%%%%%%%%%%%%%%%%%%%%%%%%%%%%%%%%%%
\subsubsection{Assuming Every Matroid Is Representable}

Many matroids cannot be represented by vectors over a given field.

Some matroids are not representable over any field.

%%%%%%%%%%%%%%%%%%%%%%%%%%%%%%%%%%%%%%%%%%%%%%%%%%%%%%%%%%%%
\subsubsection{Confusing Deletion and Contraction}

Deletion removes an element.

Contraction removes the element after treating it as already included
in the rank structure.

They are dual operations but generally produce different matroids.

%%%%%%%%%%%%%%%%%%%%%%%%%%%%%%%%%%%%%%%%%%%%%%%%%%%%%%%%%%%%
\subsubsection{Applying Greedy Optimization to Arbitrary Constraints}

Greedy selection is guaranteed for matroid independence systems.

General combinatorial constraints may not satisfy matroid exchange.

%%%%%%%%%%%%%%%%%%%%%%%%%%%%%%%%%%%%%%%%%%%%%%%%%%%%%%%%%%%%
\subsection{Applications}
\label{subsec:matroid-applications}
%%%%%%%%%%%%%%%%%%%%%%%%%%%%%%%%%%%%%%%%%%%%%%%%%%%%%%%%%%%%

\begin{applicationbox}

\textbf{Graph Theory.}
Graphic matroids abstract forests, spanning trees, cycles, cuts, and
graph minors.

\medskip

\textbf{Linear Algebra.}
Vector matroids capture linear independence, bases, span, rank, and
dependence without requiring coordinates.

\medskip

\textbf{Optimization.}
Maximum-weight basis problems are solved optimally by greedy algorithms.

\medskip

\textbf{Network Design.}
Spanning trees, forest decompositions, and connectivity constraints
have natural matroid formulations.

\medskip

\textbf{Matching Theory.}
Bipartite matching can be formulated as the intersection of two
partition matroids.

\medskip

\textbf{Finite Geometry.}
Projective and affine geometries produce representable matroids over
finite fields.

\medskip

\textbf{Coding Theory.}
Linear codes produce matroids through generator and parity-check
matrices.

\medskip

\textbf{Design Theory.}
Finite incidence geometries can simultaneously define block designs
and matroids.

\medskip

\textbf{Combinatorial Optimization.}
Matroid intersection, matroid union, and submodular optimization form
major generalizations of classical graph algorithms.

\medskip

\textbf{Algebraic Combinatorics.}
The Tutte polynomial, geometric lattices, characteristic polynomials,
and representability link matroids with algebraic structure.

\end{applicationbox}

%%%%%%%%%%%%%%%%%%%%%%%%%%%%%%%%%%%%%%%%%%%%%%%%%%%%%%%%%%%%
\subsection{Exercises}
\label{subsec:matroid-exercises}
%%%%%%%%%%%%%%%%%%%%%%%%%%%%%%%%%%%%%%%%%%%%%%%%%%%%%%%%%%%%

\begin{exercise}[1]
State the three independent-set axioms of a matroid.
\end{exercise}

\begin{exercise}[1]
Define the ground set of a matroid.
\end{exercise}

\begin{exercise}[1]
Define a basis.
\end{exercise}

\begin{exercise}[1]
Define a circuit.
\end{exercise}

\begin{exercise}[1]
Define the rank of a matroid.
\end{exercise}

\begin{exercise}[1]
What are the independent sets of a graphic matroid?
\end{exercise}

\begin{exercise}[1]
What are the circuits of a graphic matroid?
\end{exercise}

\begin{exercise}[1]
What are the bases of the graphic matroid of a connected graph?
\end{exercise}

\begin{exercise}[1]
Define the uniform matroid

\[
U_{r,n}.
\]
\end{exercise}

\begin{exercise}[1]
How many bases does

\[
U_{3,7}
\]

have?
\end{exercise}

\begin{exercise}[1]
What are the circuits of

\[
U_{3,7}?
\]
\end{exercise}

\begin{exercise}[1]
Define a loop in a matroid.
\end{exercise}

\begin{exercise}[1]
Define a coloop.
\end{exercise}

\begin{exercise}[2]
Prove that all bases of a matroid have equal cardinality.
\end{exercise}

\begin{exercise}[2]
Show that every independent set is contained in a basis.
\end{exercise}

\begin{exercise}[2]
For the graphic matroid of \(C_5\), determine:
\begin{itemize}
    \item the rank;
    \item all circuits;
    \item the number of bases.
\end{itemize}
\end{exercise}

\begin{exercise}[2]
Let \(G\) be a connected graph with \(8\) vertices. What is the rank of
its graphic matroid?
\end{exercise}

\begin{exercise}[2]
Let \(G\) have \(10\) vertices and \(3\) connected components. Find the
rank of \(M(G)\).
\end{exercise}

\begin{exercise}[2]
Show that every bridge of a connected graph is a coloop of its graphic
matroid.
\end{exercise}

\begin{exercise}[2]
Construct the vector matroid represented by

\[
A=
\begin{pmatrix}
1&0&1&1\\
0&1&1&0
\end{pmatrix}
\]

over \(\mathbb{R}\).
\end{exercise}

\begin{exercise}[2]
Determine the circuits of the vector matroid in the previous exercise.
\end{exercise}

\begin{exercise}[2]
Determine all bases of the vector matroid in the previous exercise.
\end{exercise}

\begin{exercise}[2]
Find the rank of

\[
U_{4,10}.
\]
\end{exercise}

\begin{exercise}[2]
Find the rank of a partition matroid with capacities

\[
2,3,1,4
\]

assuming each part is sufficiently large.
\end{exercise}

\begin{exercise}[2]
Give an example of a hereditary independence system that is not a
matroid.
\end{exercise}

\begin{exercise}[2]
For a matroid of rank \(r\) on \(n\) elements, determine the rank of
its dual.
\end{exercise}

\begin{exercise}[3]
Prove the basis-exchange theorem from the independent-set axioms.
\end{exercise}

\begin{exercise}[3]
Show that an independent set contains no circuit.
\end{exercise}

\begin{exercise}[3]
Show that every dependent set contains a circuit.
\end{exercise}

\begin{exercise}[3]
Prove the circuit elimination property.
\end{exercise}

\begin{exercise}[3]
Prove that the rank function is monotone.
\end{exercise}

\begin{exercise}[3]
Prove the submodularity inequality

\[
r(A\cup B)+r(A\cap B)
\leq
r(A)+r(B).
\]
\end{exercise}

\begin{exercise}[3]
For a graphic matroid, prove

\[
r(A)=|V|-c(A).
\]
\end{exercise}

\begin{exercise}[3]
Describe the closure of an edge set in a graphic matroid.
\end{exercise}

\begin{exercise}[3]
Describe the closure operation in a vector matroid.
\end{exercise}

\begin{exercise}[3]
Show that every basis is both independent and spanning.
\end{exercise}

\begin{exercise}[3]
Prove the dual rank formula

\[
r_{M^*}(A)
=
|A|-r_M(E)+r_M(E\setminus A).
\]
\end{exercise}

\begin{exercise}[3]
For a planar graph, investigate the relation

\[
M(G)^*
\cong
M(G^*).
\]
\end{exercise}

\begin{exercise}[3]
Show that deletion and contraction are exchanged by duality.
\end{exercise}

\begin{exercise}[3]
Represent the Fano matroid over

\[
\mathbb{F}_2.
\]
\end{exercise}

\begin{exercise}[3]
Determine the seven three-element circuits of the Fano matroid.
\end{exercise}

\begin{exercise}[3]
Use the greedy algorithm to find a maximum-weight basis of a partition
matroid of your own construction.
\end{exercise}

\begin{exercise}[3]
Explain why Kruskal's algorithm is an instance of matroid greedy
optimization.
\end{exercise}

\begin{exercise}[3]
Formulate bipartite matching as an intersection of two partition
matroids.
\end{exercise}

\begin{exercise}[3]
Compute the Tutte polynomial of

\[
U_{1,2}.
\]
\end{exercise}

\begin{exercise}[3]
Verify that

\[
T_M(1,1)
\]

counts the bases for a small matroid.
\end{exercise}

\begin{exercise}[4]
Study alternative matroid axiomatizations using:
\begin{itemize}
    \item bases;
    \item circuits;
    \item rank functions;
    \item closure operators.
\end{itemize}
\end{exercise}

\begin{exercise}[4]
Prove the greedy characterization theorem for matroids.
\end{exercise}

\begin{exercise}[4]
Study the matroid intersection theorem and derive the min--max formula

\[
\max |I|
=
\min_{A\subseteq E}
\left(
r_1(A)+r_2(E\setminus A)
\right).
\]
\end{exercise}

\begin{exercise}[4]
Investigate Edmonds' matroid intersection algorithm.
\end{exercise}

\begin{exercise}[4]
Study the matroid union theorem.
\end{exercise}

\begin{exercise}[4]
Investigate transversal matroids and their relationship with Hall's
Marriage Theorem.
\end{exercise}

\begin{exercise}[4]
Study gammoids and their relationship with directed linkage problems.
\end{exercise}

\begin{exercise}[4]
Investigate excluded-minor characterizations of binary matroids.
\end{exercise}

\begin{exercise}[4]
Study regular matroids and their relationship with graphic and cographic
matroids.
\end{exercise}

\begin{exercise}[4]
Investigate the Tutte polynomial of a matroid and its important
specializations.
\end{exercise}

\begin{exercise}[4]
Study the lattice of flats and its Möbius function.
\end{exercise}

\begin{exercise}[4]
Investigate the characteristic polynomial of a matroid.
\end{exercise}

\begin{exercise}[4]
Study matroid base polytopes and their role in combinatorial
optimization.
\end{exercise}

%%%%%%%%%%%%%%%%%%%%%%%%%%%%%%%%%%%%%%%%%%%%%%%%%%%%%%%%%%%%
\subsection{Conceptual Summary}
\label{subsec:matroid-summary}
%%%%%%%%%%%%%%%%%%%%%%%%%%%%%%%%%%%%%%%%%%%%%%%%%%%%%%%%%%%%

A matroid is an abstract independence structure

\[
\boxed{
M=(E,\mathcal{I})
}
\]

satisfying hereditary and exchange properties.

A basis is a maximal independent set.

All bases have the same size.

Thus the rank is well-defined:

\[
\boxed{
r(M)=|B|.
}
\]

For a subset \(A\),

\[
\boxed{
r(A)
=
\max
\{
|I|:
I\subseteq A,\,
I\in\mathcal{I}
\}.
}
\]

A circuit is a minimal dependent set.

Graphic matroids satisfy

\[
\boxed{
\text{independent sets}
=
\text{forests},
}
\]

\[
\boxed{
\text{bases}
=
\text{spanning forests},
}
\]

and

\[
\boxed{
\text{circuits}
=
\text{cycles}.
}
\]

Vector matroids abstract linear independence.

Uniform matroids satisfy

\[
\boxed{
I\text{ independent}
\iff
|I|\leq r.
}
\]

Partition matroids enforce category capacities.

Closure is defined by

\[
\boxed{
\operatorname{cl}(A)
=
\{
e:
r(A\cup\{e\})=r(A)
\}.
}
\]

Flats are closed sets.

Matroid duality defines

\[
\boxed{
\mathcal{B}(M^*)
=
\{
E\setminus B:
B\in\mathcal{B}(M)
\}.
}
\]

Deletion and contraction produce minors.

Representable matroids arise from matrices over fields.

The Fano matroid

\[
\boxed{
F_7
}
\]

comes from the Fano plane and is representable over

\[
\mathbb{F}_2.
\]

The exchange axiom gives the fundamental algorithmic property:

\[
\boxed{
\text{greedy optimization is always correct for matroid bases}.
}
\]

The Tutte polynomial is

\[
\boxed{
T_M(x,y)
=
\sum_{A\subseteq E}
(x-1)^{r(E)-r(A)}
(y-1)^{|A|-r(A)}.
}
\]

Matroid theory therefore unifies

\[
\boxed{
\text{graph cycles}
+
\text{linear independence}
+
\text{finite geometry}
+
\text{greedy optimization}
+
\text{abstract dependence}.
}
\]

%%%%%%%%%%%%%%%%%%%%%%%%%%%%%%%%%%%%%%%%%%%%%%%%%%%%%%%%%%%%
\subsection{Section Connections}
\label{subsec:matroid-connections}
%%%%%%%%%%%%%%%%%%%%%%%%%%%%%%%%%%%%%%%%%%%%%%%%%%%%%%%%%%%%

\chapterconnection{
Matroid Theory completes Chapter 6 by abstracting independence
structures that appeared throughout the chapter. Graph Theory provides
graphic matroids, where forests are independent and spanning trees are
bases. Design Theory and finite geometry provide representable matroids
such as the Fano matroid. Algebraic Combinatorics contributes the Tutte
polynomial, lattices of flats, finite-field representations, and
symmetry. Extremal and Probabilistic Combinatorics motivate independence
and optimization questions, while Basic Counting and Enumerative
Combinatorics count bases and other matroid structures. The subject also
connects strongly with Chapter~\ref{chap:algebra} through vector spaces
and fields and with Optimization and Operations Research through greedy
algorithms, matroid intersection, matroid union, and submodular
optimization.
}

%%%%%%%%%%%%%%%%%%%%%%%%%%%%%%%%%%%%%%%%%%%%%%%%%%%%%%%%%%%%
% SECTION SOURCE TRACKING
%%%%%%%%%%%%%%%%%%%%%%%%%%%%%%%%%%%%%%%%%%%%%%%%%%%%%%%%%%%%

%------------------------------------------------------------
% 6.12 Matroid Theory
%------------------------------------------------------------
%
% Source basis:
%   - Standard finite matroid theory.
%   - Standard graphic and representable matroids.
%   - Standard matroid optimization.
%   - Standard Tutte-polynomial theory.
%
% Used for:
%   - independent-set matroid axioms;
%   - hereditary property;
%   - exchange axiom;
%   - dependent sets;
%   - bases;
%   - equal basis cardinality;
%   - rank;
%   - subset rank;
%   - rank submodularity;
%   - circuits;
%   - circuit elimination;
%   - graphic matroids;
%   - graphic rank;
%   - spanning trees and forests;
%   - vector matroids;
%   - matrix representations;
%   - representable matroids;
%   - binary, ternary, and regular matroids;
%   - uniform matroids;
%   - partition matroids;
%   - loops, coloops, and parallel elements;
%   - spanning sets;
%   - closure and flats;
%   - hyperplanes;
%   - basis exchange;
%   - fundamental circuits;
%   - dual matroids;
%   - cocircuits;
%   - graphic duality;
%   - deletion;
%   - contraction;
%   - minors;
%   - excluded minors;
%   - Fano matroid;
%   - finite projective geometry matroids;
%   - greedy algorithm;
%   - greedy characterization;
%   - maximum-weight bases;
%   - matroid polytopes;
%   - matroid intersection;
%   - bipartite matching as intersection;
%   - matroid union;
%   - transversal matroids;
%   - gammoid preview;
%   - direct sums;
%   - connected matroids;
%   - restriction and truncation;
%   - Tutte polynomial;
%   - deletion--contraction;
%   - characteristic polynomial preview;
%   - lattice of flats;
%   - geometric lattices;
%   - simple and cosimple matroids;
%   - matroid isomorphisms and automorphisms;
%   - coding-theory connections;
%   - design-theory connections;
%   - optimization applications;
%   - exercises.
%
% Notes:
%   - Linear independence and field theory are developed more fully in
%     Chapter 1.
%   - Graph terminology and spanning-tree concepts were developed in
%     Section 6.9.
%   - Finite projective geometry and the Fano plane connect directly
%     with Section 6.11 Design Theory.
%   - Greedy, intersection, union, polyhedral, and submodular methods
%     connect strongly with Chapter 10 Optimization and Operations
%     Research.
%   - The Tutte polynomial was previewed in Algebraic Combinatorics and
%     Graph Theory and reaches its natural abstract setting here.
%
%%%%%%%%%%%%%%%%%%%%%%%%%%%%%%%%%%%%%%%%%%%%%%%%%%%%%%%%%%%%